\subsection{Matriz comparativa del escenario 2}

\begin{center}
\small
\begin{tabularx}{\textwidth}{|>{\raggedright\arraybackslash}p{2.6cm}|L|L|L|}
\hline
\rowcolor{verdeuteq}\textcolor{white}{\textbf{Criterio}} & \textcolor{white}{\textbf{Secuencial (cascada)}} & \textcolor{white}{\textbf{Iterativo e incremental}} & \textcolor{white}{\textbf{Ágil}} \\ \hline
K1. Estabilidad y cambios
& Si Rosa cambia de puntos a promociones después de aprobar el alcance, el equipo tendría que revisar los requisitos y el diseño afectados.
& Una primera versión de pedidos permitiría a Rosa probar el flujo y decidir en el siguiente ciclo si incluye promociones.
& Rosa puede reordenar las ideas sobre puntos y promociones al inicio de cada ciclo, según lo que más le sirva al negocio.
\\ \hline
K2. Partes interesadas y disponibilidad
& Rosa puede aprobar los requisitos al inicio y revisar el producto al final. Si sus observaciones llegan tarde, corregir el flujo puede exigir rehacer trabajo.
& Como Rosa está disponible casi a diario, puede probar cada versión y aclarar dudas antes del siguiente ciclo.
& Rosa puede decidir las prioridades con frecuencia y su hijo aportar ideas. Con dos desarrolladores a medio tiempo, centralizar esa decisión ayuda a evitar instrucciones contradictorias.
\\ \hline
K3. Tipo de sistema y criticidad
& Conviene definir desde el inicio cómo se protegen las direcciones y cómo se concilian las transferencias, porque un error puede exponer datos o dejar un pedido sin pagar.
& Una versión temprana permite probar cómo se vincula una transferencia con su pedido y corregir el flujo antes de añadir otras funciones.
& En el ciclo que incorpora transferencias, el equipo puede comprobar que un pedido solo aparezca como pagado después de conciliar el pago.
\\ \hline
K4. Riesgo principal
& Con 14 semanas y dos desarrolladores a medio tiempo, detectar tarde un fallo de conexión deja poco margen para corregirlo.
& Probar una versión temprana en un Android de gama baja con cortes de conexión permite ajustar el flujo de pedidos antes de ampliarlo.
& Con presupuesto fijo, el equipo puede priorizar menú, pedidos y transferencias. Los puntos y el pago con tarjeta pueden quedar para después.
\\ \hline
K5. Documentación y marco normativo o contractual
& Aunque el acuerdo de precio fijo sea informal, dejar por escrito el alcance y los criterios de aceptación permite revisar el impacto de los cambios.
& En cada versión se pueden registrar las decisiones de Rosa y los resultados de las pruebas de conciliación.
& En cada ciclo se pueden actualizar los requisitos sobre datos de clientes y guardar la evidencia de las pruebas de pago.
\\ \hline
\end{tabularx}
\end{center}

\subsection{Recomendación y riesgos del escenario 2}

\begin{center}
\small
\begin{tabularx}{\textwidth}{|>{\raggedright\arraybackslash}p{4.5cm}|L|}
\hline
\rowcolor{verdeuteq}\textcolor{white}{\textbf{Elemento}} & \textcolor{white}{\textbf{Respuesta}} \\ \hline
Modelo recomendado (puro o híbrido; si es híbrido, qué parte sigue cada modelo) & \\ \hline
Razón 1 (dato del escenario, criterio y consecuencia) & \\ \hline
Razón 2 & \\ \hline
Razón 3 & \\ \hline
Por qué se descartan los otros dos modelos & \\ \hline
Tensión principal del escenario (el dato que apunta a otro modelo) & \\ \hline
Riesgo 1 del modelo elegido y su mitigación & \\ \hline
Riesgo 2 del modelo elegido y su mitigación & \\ \hline
\end{tabularx}
\end{center}
